\documentclass{ximera}
\title{Announcements, HW, Math 151 Online}

\newcommand{\pskip}{\vskip 0.1 in}

\begin{document}
\begin{abstract}
Announcements, HW, Math 151 Online, Spring 2026
\end{abstract}
\maketitle


\section{Announcements}
\begin{enumerate}
\item {\bf HW5B:} Due to the Canvas outage, the version of HW5B below is shorter than the original. It will be due after Canvas comes back online.

\item {\bf HW6A:} (May 10) See below for HW6A. We'll see what happens with Canvas, but there is a possibility you will not be required to submit this assignment. But you will be responsible for its content.

\item {\bf Exam 2:} (May 10) Our second exam is this Friday. Please be sure to register with your testing center.

{\bf Required Materials:}
\begin{enumerate}
\item A ruler to draw straight lines and to make measurements.

\item Pencils and an eraser.

\item That's all. No other materials (ie. phones, dictionaries, calculators, notes, etc.)
\end{enumerate}

{\bf Scope:} The exam is cumulative, from HW0A to HW5B. It will focus on the chain rule and trigonometry. There will be a question about interpreting a derivative as a scaling factor that requires you give a example of how the scaling factor works with \emph{specific} small changes. See the solution to Question 4c in Exam 1 and the two problems here to be sure you know what this means:

\href{https://ximera.osu.edu/calcone/Calculus1/ScalingFactorsSummary/ScalingFactorsSummary}{The Derivative as a Scaling Factor}

{\bf Directions:}

Here are the directions for the exam:

\begin{itemize}
\item Show all work. Do not skip steps.

\item Use the Leibniz notation for all derivatives and their evaluations.

\item Do not use the quotient rule.

\item Show \emph{all} steps in using the chain rule {\bf exactly} as in our class notes. You are resonsible for knowing what this means. See the Explanations to Examples 1,2, and 3 here: \href{https://ximera.osu.edu/calcone/Calculus1/ChainRule/ChainRule}{The Chain Rule} for examples.

\item Work vertically, with at most one equal sign per line. Do not split the page into columns.

\item You may be asked to find an equation of a tangent line to a curve. See this link for an example:

\href{https://ximera.osu.edu/calcone/Calculus1/TrigTangentLine/TrigTangentLine}{Trigonometry Tangent Line}
%\end{enumerate}

\item There is a two-hour time limit. The exam ends two hours after you begin.

\item You may {\bf not} ask any questions about this test. This includes all types of questions, including questions about meaning or asking for a translation.

\item Do {\bf not} share any information about this exam with your classmates.

\end{itemize}

\pskip

\noindent {\bf Identities:}

\[
 \sin (A-B) = \sin A \cos B - \cos A \sin B
\]


\end{enumerate}


\section{Homework 5B}

\begin{enumerate}
\item Read the chapter \href{https://ximera.osu.edu/calcone/Calculus1/InverseShortest/InverseShortest}{Derivatives of Inverse Functions, Shortest Version} of our class notes. 

\begin{enumerate}
\item {\bf Required:} State and answer the Free Response Questions (those with the blue Submit button) in Questions 1, 2, and 3. Do not submit your answers in Ximera, but with your homework when Canvas reopens.

\item {\bf Required:} State and fill in the Question Boxes in Question 1 and Theorem 2.

\item {\bf Required:} Exercises 4, 5, and 6 (all parts) in this chapter.
\end{enumerate}

\item {\bf Recommended:} Questions 17, 18, 21, and 22 in the chapter \href{https://ximera.osu.edu/calcone/Calculus1/InverseF/InverseF}{Derivatives of Inverse Functions} of our class notes.

\item {\bf Recommended:} Questions 1-5 in the chapter \href{https://ximera.osu.edu/calcone/Calculus1/InverseTrigCW/InverseTrigCW}{InverseTrigCW} of our class notes. 

\end{enumerate}

\section{Homework 6A:} State and answer all the questions in this chapter of our class notes:

\href{https://ximera.osu.edu/calcone/Calculus1/RelativeChanges/RelativeChanges}{Relative Changes and Rates of Change}




\end{document}